\documentclass{article}
\usepackage[utf8]{inputenc}
\usepackage[T1]{fontenc}
\usepackage{geometry}
\geometry{
    a4paper,
    left=2.5cm,
    right=2.5cm,
    top=2.5cm,
    bottom=2.5cm
}
\usepackage{graphicx}
\usepackage{float}
\usepackage{hyperref}
\usepackage[english]{babel}
\usepackage{amsmath}
\usepackage{amssymb}
\usepackage{array}
\usepackage{booktabs}
\usepackage{caption}
\usepackage{subcaption}
\usepackage{enumitem}
\usepackage{fancyhdr}
\usepackage{titlesec}
\usepackage{url}
\usepackage{textcomp}
\usepackage{xcolor}
\usepackage{listings}
\usepackage{multirow}
\usepackage{longtable}
\usepackage{tabularx}
\usepackage{times}

% Set up page style
\pagestyle{fancy}
\fancyhf{}
\rfoot{\thepage}
\renewcommand{\headrulewidth}{0pt}

% Title formatting
\titleformat{\chapter}[display]{\normalfont\bfseries}{\chaptertitle\thechapter}{0pt}{\Huge}
\titlespacing*{\chapter}{0pt}{50pt}{40pt}

% Document info
\title{Template Document}
\author{Academic Support Agent}
\date{\today}

\begin{document}
\maketitle

% Inserindo o cabe alho
	\cabecalho
	% Inserindo as informa es da prova e garantindo que estar o centralizados na p gina.
	\info
	
	% Descomente a linha seguinte caso voc  queira gerar um PDF com as solu es  s quest es.
	\printanswers
	
	%---------------------------------------------------------------------------
	% QUEST ES -----------------------------------------------------------------
	\begin{questions} % Ambiente de quest es
		% Comando para deixar os textos justificados.   necess rio repetir o comando para cada novo ambiente.
		\justifying
		% Cada novo comando \question iniciar  uma nova quest o. A numera o ser  adicionada automaticamente.
		\question[1] % O n mero dentro dos colchetes   o valor da quest o.
		Este   o exemplo de uma quest o: o que significa o valor entre colchetes?
		%-----------------------------------------------------------------------
		\question[1\half] % N o   poss vel uma quest o ter pontua o "quebrada", ou seja, 1,5. Mas   poss vel expressar com fra es. Outra solu o   expressar a pontua o entre 0 e 100.
		Este   o exemplo de uma quest o contendo equa es ou f rmulas:\\ 
		$x^{2} + 2\pi + \pi^{2}$
		%-----------------------------------------------------------------------
		\question % Definir a pontua o   opcional. Ao se definir junto do comando \question o LaTeX faz c lculos para mostrar a quantidade total de pontos da prova. N o h  necessidade desse c lculo, portanto, caso o professor queira, talvez seja melhor acrescentar a pontua o da forma que melhor lhe convir.
		Mais uma quest o de exemplo: a defini o da pontua o   opcional [1,0].
		%-----------------------------------------------------------------------
		\question
		Uma quest o pode tamb m ser dividida em partes listadas como letras.
			% Come ando o ambiente de partes
			\begin{parts}
				\justifying % Tem de ser declarado em todo novo ambiente
				\part[1] Essa   a primeira letra. % Uma parte pode tamb m ter pontua o
				\part 
				Fazendo essa quest o encontrar  a resposta? % Pode ficar na linha seguinte, que nem as quest es
				\part Na verdade a resposta voc  encontra ao resolver a letra (\ref{letraE}).
				\part Resolva a (\ref{letraE}).
				\part Aqui voc  encontra a resposta.\label{letraE}
					\begin{subparts}
						\subpart Ou ser  que n o? % Subpartes, e subsubpartes s o permitidos.
					\end{subparts}
			\end{parts}
		%-----------------------------------------------------------------------
		\question
		Deixando um espa o para que os alunos possam resolver a quest o.
		\vspace{3cm} % v[ertical]space
		% Caso voc  queira distribuir os espa os igualmente enter todas as quest es, e partes, subpartes e subsubpartes, basta escrever o seguinte comando ao fim da quest o:
		%			\vspace{\stretch{1}}
		% Se alguma quest o necessita do dobro de espa o em rela o  s demais, basta escrever ao fim da quest o o seguinte comando:
		%			\vspace{\stretch{2}}
		% O mesmo racioc nio se aplica ao triplo do espa o, etc.
		%-----------------------------------------------------------------------
		\question
		A quest o seguinte vem depois do espa o dado para a quest o anterior.
			\begin{parts}
				\part O espa amento pode ser usado nas partes tamb m. \vspace{1cm}
				\part A parte seguinte vem depois de um espa o.
			\end{parts}
		%-----------------------------------------------------------------------
		% Caso queira mudar logo de p gina, garantindo o espa o restante para a  ltima quest o, basta utilizar o comando a seguir:
		\newpage
		\question
		Quest o em uma nova p gina. E d  para deixar linhas para os alunos responderem.
		\fillwithlines{5cm}
		%-----------------------------------------------------------------------
		\question
		Uma quest o de m ltipla escolha:
			\begin{choices}
				\choice Essa   a alternativa correta.
				\choice Na verdade   essa.
				\choice N o   a anterior.
				\choice Mas talvez seja.
				\choice Melhor marcar essa aqui.
			\end{choices}
		%-----------------------------------------------------------------------
		\question
		Se as alternativas puderem ser em somente uma linha, basta utilizar o ambiente \texttt{oneparchoices}.
			\begin{oneparchoices}
				\choice Eu
				\choice Ele
				\choice Ela
				\choice N s
				\choice Ningu m
			\end{oneparchoices}
		%-----------------------------------------------------------------------
		\question
		Mas tamb m   poss vel fazer assim:
			\begin{enumerate}[a)]
				\item Primeira alternativa.
				\item Segunda alternativa.
				\item Terceira alternativa.
				\item Quarta alternativa.
				\item Quinta alternativa.
			\end{enumerate}
		%-----------------------------------------------------------------------
		\question
		Utilizando \texttt{multicols} para separar as alternativas em uma mesma linha:
			\begin{multicols}{5}
				\begin{enumerate}[(A)]
					\item Alternativa 1
					\item Alternativa 2
					\item Alternativa 3
					\item Alternativa 4
					\item Alternativa 5
				\end{enumerate}
			\end{multicols}
		%-----------------------------------------------------------------------
		\question
		Uma quest o de m ltipla escolha pode ter uma alternativa j  marcada como a correta:
			\begin{choices}
				\justifying
				\CorrectChoice A alternativa correta vai aparecer se a op o \texttt{printanswers} estiver ativa no in cio do documento.
				\choice Alternativa errada.
				\choice Alternativa errada.
				\choice Alternativa errada.
				\choice Alternativa errada.
			\end{choices}
		%-----------------------------------------------------------------------
		\question
		\justifying
		Uma quest o pode ter outros tipos de solu o (no arquivo UninassauProva.sty tem um trecho onde   poss vel fazer outras configura es no ambiente de solu o):
			\begin{solution}
				\justifying Solu o utilizando o ambiente \texttt{solution}. Se \texttt{printanswers} estiver desabilitado, voc  deve fornecer o tamanho do espa o em branco que ser  posto no lugar da solu o. Caso esse valor n o seja especificado, o LaTeX n o acrescentar  espa o em branco.\\ 
                Ex.: \verb*|\begin{solution}[5cm]|
			\end{solution}
			\begin{solutionorbox}
				Solu o com o ambiente \texttt{solutionorbox}. Se a solu o n o for impressa, haver  no lugar uma caixa de texto em branco, do tamanho especificado pelo usu rio.\\ 
                Ex.: \verb*|\begin{solutionorbox}[2cm]|
			\end{solutionorbox}
			\begin{solutionorlines}
				Solu o com o ambiente \texttt{solutionorlines}. Se a solu o n o for impressa, haver  no lugar linhas para o aluno escrever.\\
                Ex.: \verb*|\begin{solutionorlines}[2cm]|
			\end{solutionorlines}
			\begin{solutionordottedlines}
				Solu o com o ambiente \texttt{solutionordottedlines}. Se a solu o n o for impressa, haver  no lugar linhas pontilhadas para o aluno escrever.\\
                Ex.: \verb*|\begin{solutionordottedlines}[2cm]|
			\end{solutionordottedlines}
			\begin{solutionorgrid}
				Solu o com o ambiente \texttt{solutionorgrid}. Se a solu o n o for impressa, haver  no lugar uma grade.\\
                Ex.: \verb*|\begin{solutionorgrid}[2cm]|
			\end{solutionorgrid}
		%-----------------------------------------------------------------------
		% A linha abaixo serve para formatar a apar ncia da quest o b nus.
		\bonusqformat{\textbf{Quest o b nus} \dotfill [\thepoints]}
		\bonusquestion[1]
		  poss vel acrescentar tamb m quest es b nus.
	\end{questions}
\end{document}